\documentclass[../RFC-HDFG-2026-001.tex]{subfiles}

\begin{document}

\section{Introduction}
\label{sec:intro}

\subsection{Purpose}

This document describes the high-level design for a new unified filter
configuration interface for HDF5 I/O filters. The feature adds a usability
layer on top of the existing filter machinery, accepting filter parameters as
either raw \texttt{cd\_values} arrays or human-readable \texttt{key=value}
strings through a single new API, \texttt{H5Pappend\_filter}.

\subsection{Scope}

This design covers:

\begin{itemize}
\item New public API functions in the \texttt{H5P} and \texttt{H5Z} modules for
  filter configuration, including \texttt{H5Pappend\_filter} and the
  \texttt{H5Z\_params\_t} tagged union type.
\item A new version of the filter plugin class (\texttt{H5Z\_class3\_t}) with optional
  string-configuration callbacks.
\item An on-disk format change: a new version~3 of the filter-pipeline
  object-header message (\texttt{H5O\_PLINE\_VERSION\_3},
  \SectionRef{sec:pline-v3}) that persists the validated parameter string
  verbatim alongside \texttt{cd\_values}, gated by
  \texttt{H5Pset\_libver\_bounds}.
\item Integration points for HDF5 command-line tools (\texttt{h5repack}, \texttt{h5dump}).
\item The Fortran and Java bindings that ship with HDF5
  (\SectionRef{sec:bindings}): a generic
  \texttt{h5pappend\_filter\_f} with two specific procedures plus
  \texttt{h5pget\_filter\_params\_by\_idx\_f} and
  \texttt{h5zconfig\_get\_param\_f} for Fortran; two overloaded
  \texttt{H5Pappend\_filter} methods and
  \texttt{H5Pget\_filter\_params\_by\_idx} for Java.
\end{itemize}

This design does not cover:

\begin{itemize}
\item Changes to the filter execution pipeline itself.
\item External high-level language binding implementations (h5py, etc.), though
  the design is motivated in part by their needs.
\item Data-attribute-dependent compression strategies (e.g., choosing different
  filters based on dimensionality or datatype), which could be layered on top
  of this feature in the future.
\item A string-based equivalent of \texttt{H5Pmodify\_filter}. Users who need
  to update the parameters of an existing filter entry in a copied DCPL should
  retrieve the current \texttt{cd\_values} via \texttt{H5Pget\_filter\_by\_id2}
  and call \texttt{H5Pmodify\_filter} directly.
\item VFD and VOL connector string configuration.  The TOML-based parameter
  string grammar defined here is designed to extend naturally to VFD/VOL
  configuration, and \SectionRef{sec:vfd-use-cases} documents that design
  as forward-looking context.  The VFD/VOL configuration API changes
  (\texttt{H5FDconfig\_get\_subtable\_str}, profile constants, etc.) are
  explicitly out of scope for this RFC.
\item In-file blob configuration storage for filters whose configuration
  data is too large for \texttt{cd\_values} or the parameter string (e.g.,
  pre-processed source for a JIT-compiled filter module, or a reference to
  a spatial mask stored as a separate dataset).  This is addressed by a
  follow-on RFC, RFC-HDFG-2026-003, summarized as forward-looking context
  in \SectionRef{sec:blob-future}.
\end{itemize}


\section{Motivation}
\label{sec:motivation}

\subsection{Problem Statement}

The current generic interface for configuring HDF5 filters is \texttt{H5Pset\_filter},
which accepts an array of \texttt{unsigned int} values (\texttt{cd\_values}). This interface
presents three categories of problems.

\verysubsection{Type punning}
Filters that accept floating-point parameters (rates,
tolerances, precision thresholds) require users to manually bit-cast \texttt{double}
or \texttt{float} values into one or more \texttt{unsigned int} slots. This is error-prone,
endian-sensitive, and difficult or impossible to express naturally in
high-level languages such as Python, Java, and MATLAB without dedicated
wrapper code. For example, configuring H5Z-ZFP in rate mode currently
requires:

\begin{minted}{c}
unsigned int cd_vals[4];
double rate = 3.5;
cd_vals[0] = H5Z_ZFP_MODE_RATE;
memcpy(&cd_vals[2], &rate, sizeof(double));
H5Pset_filter(plist, H5Z_FILTER_ZFP, flags, 4, cd_vals);
\end{minted}

\verysubsection{CLI usability}
Command-line tools such as \texttt{h5repack} require users to pass
opaque integer sequences (e.g., \texttt{-f UD=32013,\allowbreak{}0,\allowbreak{}4,\allowbreak{}1,\allowbreak{}0,\allowbreak{}0,\allowbreak{}1074921472}). External
helper scripts exist solely to translate human-readable settings into these
sequences.

\verysubsection{Opaque introspection}
When examining a file's dataset creation properties,
the \texttt{cd\_values} array is meaningless without external documentation. There is
no standard way to recover the original intent of the configuration (e.g.,
``ZFP rate mode, rate=3.5'').

\subsection{Goals}

\begin{enumerate}
\item Users can configure any filter using human-readable named
   parameters.
\item Configured filter pipelines are introspectable in human-readable form,
   with both the filter's label and its configuration string persisting in
   the file independent of plugin availability, and round-tripping without
   loss.
\item All existing APIs, plugins, and call sites continue to work without
   modification.
\item The design supports both built-in and dynamically loaded plugin
   filters.
\item The API is consistent with the structure and conventions of existing
   HDF5 public APIs.
\end{enumerate}

\end{document}
